\documentclass[journal]{IEEEtran}
% Alternative options:
% \documentclass[conference]{IEEEtran}  % For conference papers
% \documentclass[letterpaper]{IEEEtran}  % For US letter size

% ============================================================================
% PACKAGES
% ============================================================================

\usepackage{amsmath,amsfonts}      % Mathematical symbols and environments
\usepackage{algorithmic}            % Algorithm pseudocode
\usepackage{algorithm}              % Algorithm float environment
\usepackage{graphicx}               % Image insertion
\usepackage{cite}                   % Citation management
\usepackage{booktabs}               % Professional tables
\usepackage[unicode]{hyperref}      % Hyperlinks

% Chinese support (remove for English-only papers)
\usepackage{ctex}

% Image path
\graphicspath{{figures/}}

% Custom hyphenation (add more as needed)
\hyphenation{op-tical net-works semi-conduc-tor IEEE-Xplore}

% ============================================================================
% CUSTOM COMMANDS FOR WIRELESS COMMUNICATIONS
% ============================================================================

% Mathematical symbols
\newcommand{\E}{\mathbb{E}}           % Expectation
\newcommand{\R}{\mathbb{R}}           % Real numbers
\newcommand{\C}{\mathbb{C}}           % Complex numbers
\newcommand{\bx}{\mathbf{x}}          % Bold vectors
\newcommand{\by}{\mathbf{y}}          % Bold vectors
\newcommand{\bz}{\mathbf{z}}          % Bold vectors
\newcommand{\bw}{\mathbf{w}}          % Bold vectors
\newcommand{\bH}{\mathbf{H}}          % Bold matrices
\newcommand{\bG}{\mathbf{G}}          % Bold matrices
\newcommand{\bW}{\mathbf{W}}          % Bold matrices
\newcommand{\bI}{\mathbf{I}}          % Identity matrix

% Functions
\newcommand{\norm}[1]{\left\lVert#1\right\rVert}
\newcommand{\abs}[1]{\left\lvert#1\right\rvert}

% ============================================================================
% DOCUMENT BEGIN
% ============================================================================

\begin{document}

% ----------------------------------------------------------------------------
% TITLE AND AUTHORS
% ----------------------------------------------------------------------------

\title{Resource Allocation for NOMA Systems\\with Deep Reinforcement Learning}

\author{Author Name\thanks{Manuscript received xxxx, 2024; revised xxxx, 2024. This work was supported by the National Science Foundation under Grant XXXXXXXX.}% <-this % stops a space
\thanks{The authors are with the Department of Electrical Engineering, University Name, City, Country (e-mail: author@university.edu).}
}

% The paper headers
\markboth{IEEE Transactions on Communications,~Vol.~xx, No.~xx, xxxx,~2024}%
{Author Name \MakeLowercase{\textit{et al.}}: Resource Allocation for NOMA}

% Copyright notice (optional)
%\IEEEpubid{0000--0000/00\$00.00~\copyright~2024 IEEE}

\maketitle

% ----------------------------------------------------------------------------
% ABSTRACT
% ----------------------------------------------------------------------------

\begin{abstract}
Non-orthogonal multiple access (NOMA) has emerged as a promising technique for
enhancing spectral efficiency in 5G networks. However, conventional power
allocation schemes fail to adapt to dynamic channel conditions and user
requirements. This paper proposes a deep reinforcement learning-based power
allocation framework that optimizes sum-rate while maintaining user fairness.
Our approach employs a multi-agent deep deterministic policy gradient algorithm
that learns optimal power allocation policies through continuous interaction
with the wireless environment. Simulation results demonstrate that the proposed
scheme achieves 23\% higher sum-rate and 35\% improvement in user fairness
compared to conventional fixed power allocation.
\end{abstract}

\begin{IEEEkeywords}
NOMA, deep reinforcement learning, power allocation, 5G, user fairness
\end{IEEEkeywords}

% ----------------------------------------------------------------------------
% INTRODUCTION
% ----------------------------------------------------------------------------

\section{Introduction}
\label{sec:introduction}

\IEEEPARstart{T}{he} exponential growth of mobile data traffic has driven the
development of fifth-generation (5G) wireless networks. Among the key enabling
technologies for 5G, non-orthogonal multiple access (NOMA) has attracted
significant attention due to its potential to improve spectral efficiency.

Despite the theoretical benefits of NOMA, practical implementation faces several
challenges. Conventional power allocation schemes typically rely on fixed power
ratios determined by channel state information, which fail to adapt to dynamic
network conditions and varying user requirements.

Several approaches have been proposed to address power allocation in NOMA
systems. Game-theoretic methods provide distributed solutions but often
converge to local optima. Optimization-based approaches guarantee optimality
but require complete channel state information and high computational
complexity. Recent works have explored machine learning techniques, yet most
rely on supervised learning requiring extensive labeled training data.

This paper proposes a deep reinforcement learning-based power allocation
framework for NOMA systems that addresses the limitations of existing approaches.
Our main contributions are:
\begin{itemize}
    \item We formulate the power allocation problem as a Markov decision process.
    \item We design a multi-agent deep deterministic policy gradient algorithm.
    \item We demonstrate significant improvements in sum-rate and user fairness.
\end{itemize}

The remainder of this paper is organized as follows. Section~\ref{sec:system_model}
presents the system model. Section~\ref{sec:proposed_method} describes the
proposed method. Section~\ref{sec:results} presents simulation results.
Section~\ref{sec:conclusion} concludes the paper.

% ----------------------------------------------------------------------------
% SYSTEM MODEL
% ----------------------------------------------------------------------------

\section{System Model}
\label{sec:system_model}

Consider a downlink single-cell NOMA system where a base station (BS) equipped
with $N_t$ antennas serves $K$ single-antenna users. The channel vector between
the BS and user $k$ is denoted as $\mathbf{h}_k \in \mathbb{C}^{N_t \times 1}$.

The transmitted signal at the BS is given by
\begin{equation}
\label{eq:transmitted_signal}
\mathbf{x} = \sum_{k=1}^K \sqrt{p_k}s_k,
\end{equation}
where $p_k$ is the allocated power for user $k$ and $s_k$ is the normalized data
symbol with $\mathbb{E}[|s_k|^2] = 1$. The received signal at user $k$ is
\begin{equation}
\label{eq:received_signal}
y_k = \mathbf{h}_k^H\mathbf{x} + n_k,
\end{equation}
where $n_k \sim \mathcal{CN}(0, \sigma^2)$ is additive white Gaussian noise.

% ----------------------------------------------------------------------------
% PROPOSED METHOD
% ----------------------------------------------------------------------------

\section{Proposed Method}
\label{sec:proposed_method}

\subsection{Problem Formulation}

The power allocation problem is formulated as
\begin{align}
\label{eq:optimization}
\max_{p_1,\ldots,p_K} \quad & \sum_{k=1}^K R_k \\
\text{s.t.} \quad & \sum_{k=1}^K p_k \leq P_{\mathrm{total}}.
\end{align}

\subsection{Deep Reinforcement Learning Framework}

We employ a deep deterministic policy gradient (DDPG) algorithm to learn the
optimal power allocation policy.

% ----------------------------------------------------------------------------
% SIMULATION RESULTS
% ----------------------------------------------------------------------------

\section{Simulation Results}
\label{sec:results}

\subsection{Simulation Setup}

We consider a single-cell NOMA system with $K=10$ users. Table~\ref{tab:parameters}
summarizes the simulation parameters.

\begin{table}[htbp]
\centering
\caption{Simulation Parameters}
\label{tab:parameters}
\begin{tabular}{lc}
\hline
Parameter & Value \\
\hline
Number of users ($K$) & 10 \\
Number of BS antennas ($N_t$) & 4 \\
Cell radius & 500 m \\
Path loss exponent ($n$) & 3.7 \\
Transmit power ($P_{\mathrm{total}}$) & 43 dBm \\
Noise power ($\sigma^2$) & -90 dBm \\
\hline
\end{tabular}
\end{table}

\subsection{Performance Comparison}

Figure~\ref{fig:results} shows the performance comparison. The proposed scheme
achieves significant improvements over baseline methods.

\begin{figure}[htbp]
\centering
% Replace with actual figure: \includegraphics[width=0.8\columnwidth]{figures/results.pdf}
\fbox{\parbox{0.8\columnwidth}{\centering\vspace{2cm}[Insert Figure Here]\vspace{2cm}}}
\caption{Performance comparison of different schemes.}
\label{fig:results}
\end{figure}

% ----------------------------------------------------------------------------
% CONCLUSION
% ----------------------------------------------------------------------------

\section{Conclusion}
\label{sec:conclusion}

This paper proposed a deep reinforcement learning-based power allocation
framework for NOMA systems. Simulation results demonstrate that the proposed
scheme achieves significant improvements in both sum-rate and user fairness
compared to conventional schemes. Future work will extend the framework to
multi-cell scenarios.

% ----------------------------------------------------------------------------
% ACKNOWLEDGMENT
% ----------------------------------------------------------------------------

\section*{Acknowledgment}
The authors would like to thank the anonymous reviewers for their valuable
comments and suggestions.

% ----------------------------------------------------------------------------
% REFERENCES
% ----------------------------------------------------------------------------

\bibliographystyle{IEEEtran}
\bibliography{myref}

% ============================================================================
% END DOCUMENT
% ============================================================================

\end{document}
